\section{Introduction}
\label{sec:introduction}

Surface electromyography (EMG) pattern recognition is a widely used approach to multifunction myoelectric control \citep{hudgins1993new,englehart2003robust,scheme2011electromyogram,farina2014extraction}. Recognition generally improves with more recording sites, but every electrode adds cost, wiring, contact area and a point of failure in a wearable device \citep{vujaklija2016new}. This makes it attractive to identify a reduced set of electrodes that preserves recognition performance \citep{phinyomark2012feature,liu2014emg,altimemy2013classification,stango2015spatial,qu2021reduce}.

Three methodological problems limit what published electrode-reduction results can tell a device designer. First, most reductions are selected and evaluated within the same users. User-dependent evaluation is appropriate for personalised systems, but a device intended for many users faces a different question: a layout chosen for one person need not suit the next, and subject-independent recognition is generally harder than user-dependent recognition \citep{phinyomark2018emg,cote2019deep}. In addition, if data from a test user influence scaling, ranking or tuning, the reported subset is optimistically biased. Second, feature-attribution methods such as Shapley additive explanations \citep{lundberg2017unified,lundberg2020local} are increasingly used to rank channels \citep{salmi2025shapaad}, but attributions can be unstable across data splits \citep{visani2022statistical,kelodjou2023shaping,xiang2026roshap}, and a ranking computed on the evaluation data is itself a form of leakage. Third, the number of electrodes to retain is usually read off an accuracy curve by eye, with no prespecified tolerance and no statement of what `no meaningful loss' means \citep{schuirmann1987comparison,walker2011understanding}. A margin chosen after the curve is seen cannot support a non-inferiority claim \citep{campbell2018noninferiority}.

This study tests explainability-guided electrode reduction under conditions designed to remove these three sources of optimistic bias. Each held-out participant is excluded from feature standardisation, model fitting and electrode ranking. Rankings are computed by cross-fitting: TreeSHAP attributions are always evaluated on participants whom the attributing model did not see. The primary decision rules were fixed in a statistical analysis plan before any result with a reduced electrode set was available: the non-inferiority margin ($\delta = 0.02$ in Macro-F1), the participant as the statistical unit, the bootstrap, Holm correction within each confirmatory family of tests, and the definition of the smallest acceptable electrode count $k^*$. The detailed tests for the secondary hypotheses were specified later, and all departures from the plan are reported. We also test a stability-regularised cross-fitted attribution score (SR-CFA), which penalises electrodes whose attribution varies across disjoint groups of training participants.

The study is primarily a test, not a method proposal. Using the 43-participant GRABMyo dataset \citep{pradhan2022multiday}, we asked four prespecified questions:
(H1) is some subset of $k < 16$ electrodes non-inferior to the full array?
(H2) do attribution-ranked subsets outperform random subsets of the same size?
(H3) does the stability penalty improve on attribution magnitude alone?
(H4) is the attribution faithful, that is, does removing highly ranked electrodes reduce performance more than removing random or low-ranked ones?
We report every answer, including the negative ones. Exploratory analyses covered transfer across recording sessions, bipolar re-referencing of the montage, and methodological replication on NinaPro DB5 \citep{pizzolato2017comparison}. A post-hoc extension tests whether estimating attribution heterogeneity across individual participants changes electrode selection. Figure~\ref{fig:overview} gives an overview of the study.

\begin{figure}[t]
\centering
\includegraphics[width=\textwidth]{figures/fig1_overview.pdf}
\caption{Study overview. (A) Documented forearm electrode topology, shown unwrapped; the ends of each row meet around the forearm. Dotted lines link electrodes used to derive bipolar pairs. The schematic makes no channel-to-muscle assignment. (B) Original fold-level aggregation and the post-hoc participant-level extension. In both cases, attributions come only from participants held out of the inner attributing model and exclude the outer test participant. No performance result is illustrated.}
\label{fig:overview}
\end{figure}
